\documentclass[../main2.tex]{subfiles}
\begin{document}

\begin{motivation}
2008年秋天，华尔街拥有全世界最多的数据、最贵的博士、最复杂的模型。然后它预测错了一切。\\
\textbf{核心动机}：这本书不教你预测股价，它问一个更尖的问题——为什么一群最会算的人，预测得最离谱？答案不在数学，在“因果”。本书要给的，是一套从法庭到科学界都在用的、拆解“谁真正推动了结果”的工具。
\end{motivation}

\begin{logicmap}
\begin{tikzpicture}[node distance=0.9cm, every node/.style={draw, rectangle, rounded corners, align=center, font=\small}]
\node (q) {钩子：算得最准的人，预测最离谱};
\node (r) [below=of q] {诊断：只看过去，没看动机与因果};
\node (s) [below=of r] {承诺：一套拆解因果的工具};
\node (h) [below=of s] {全书地图：动机→因果→法庭→齿轮→循环};
\node (c) [below=of h] {检验标准：一句话立论};
\draw[-{Stealth}] (q) -- (r);
\draw[-{Stealth}] (r) -- (s);
\draw[-{Stealth}] (s) -- (h);
\draw[-{Stealth}] (h) -- (c);
\end{tikzpicture}
\end{logicmap}

% ==================================================================
\section{钩子：为什么算得最准的人，预测得最离谱}

\begin{readingnote}
你有没有想过，2008年亏得最惨的，往往不是啥都不懂的小散户，而是数据最多、博士最多的那拨人？\\
多数人以为信息越多越准，其实信息越多，越容易把“描述过去”和“预测未来”搞混。
\end{readingnote}

2007年夏，贝尔斯登的两只基金先倒下。2008年9月，雷曼兄弟倒下。那几周，华尔街的风险模型每天吐出几十页报告，数字精确到小数点后四位。模型说，房价同时下跌的概率极低。结果房价真的同时跌了。

但是，为什么模型算得那么细，却错得那么彻底？

因为模型在回答一个问题，世界在问另一个问题。模型答的是“过去发生过什么”。世界问的是“接下来会发生什么”。这是两件事。

过去十年，房价一直在涨，违约率很低。模型记住了这个相关性。它把“过去相关”当成了“未来因果”。一旦环境变了，相关性断了，它连自己之前为什么准都解释不了。

这不是华尔街的专利。苏联崩溃前，情报机构年年写“体制稳定”。评级机构把垃圾债贴上AAA。共同点不是数据少，是把“能描述”当成了“能预测”。

\begin{questionbox}
\textit{现在你可能想反驳：可他们模型在平稳年代明明很准啊，你不能说他们完全不会算。}\\
\textbf{准，是因为那段日子没出意外。}平稳年代，相关性看起来像因果，拟合过去就够了。暴风雨来之前，你觉得自己看懂了天气。暴风雨一来，你才发现自己只会背过去的天气表。这是最危险的一种不会算。
\end{questionbox}

\paragraph*{逻辑衔接}
我们看到失败的表象：算得最准，预测最离谱。但失败的机制是什么？不是缺数据，是缺一种更古老的能力：看懂局中人的动机，以及分清“两件事相关”和“一件事导致了另一件”。下一节“诊断：他们懂统计，却不懂因果”将拆开这个机制。

% ==================================================================
\section{诊断：他们懂统计，却不懂因果}

\begin{readingnote}
你有没有遇过这种人：能背出一堆相关数字，却认真建议“夏天少卖冰淇淋来防淹死”？\\
多数人以为相关性就是原因，其实相关性只说“它们一起出现”，没说“谁推了谁”。
\end{readingnote}

冰淇淋销量和溺水人数正相关。夏天热，两者都涨。把冰淇淋当成溺水的原因，就是把相关当因果。

华尔街的模型也一样。它记录“房价涨，违约低”，没记录“为什么违约低”。房价涨的时候，大家都能转手，违约自然低。房价不涨了，违约就上来了。模型没写进“人为什么要还钱”“人什么时候还不起”——也就是动机。

但是，为什么动机这么重要？

因为社会不是天气。天气不关心你怎么想，社会全由人怎么想驱动。房价为什么涨？因为有人在买，有人在做空，有人在被KPI逼着包装风险。只看价格曲线，你看到的是结果。看局中人的动机，你才看到驱动。

这里要先交个底：能预测，和有因果，其实是一回事。你能预测“改变A之后的结果”，说明你知道A是不是原因。你知道A是不是原因，你就能预测“改变A之后的结果”。本书把目标限定得很窄——\textbf{预测“社会在不同选择下会走向哪里”}，不是预测趋势。前者需要因果，后者只需要拟合。

\begin{questionbox}
\textit{现在你可能想反驳：天气预报不也预测吗？它管什么因果？}\\
\textbf{问得好，这正是本书要绕开的坑。}天气预报预测的是趋势：给定现在的大气状态，推下一步。它答不了“如果我人工降雨，会不会下雨”。而本书要干的，恰恰是这种“如果当初没有X，今天会怎样”的干预性预测。没有因果，你答不了这个问题。
\end{questionbox}

\begin{reader_anticipation}
\item 天气预报不需要因果也能准，社会预测为什么不行？→ 因为本书要预测的是“如果我们换个选择会怎样”，这是干预性预测，Part II会把“趋势预测”和“干预预测”分清楚。
\item 动机看不见摸不着，怎么放进预测？→ Part I会把“想要什么”写成可比较的排序，再从行为里反推。第1章从一句最日常的话开始：“他这么做，是因为……”。
\end{reader_anticipation}

\paragraph*{逻辑衔接}
诊断完了：缺的是因果，而因果的根基是动机。光说“要有因果”还是空话，工具长什么样？下一节“承诺：一套从法庭到科学界都在用的工具”给出全书的承诺：不给水晶球，给一把拆解因果的螺丝刀。

% ==================================================================
\section{承诺：一套从法庭到科学界都在用的工具}

\begin{readingnote}
你有没有想过，为什么法官必须说“70\%是他的责任”，而经济学家可以说“大概有关”？\\
多数人以为法庭是讲道德的地方，其实法庭是人类最早被迫把因果算成数字的地方。
\end{readingnote}

同样要“分锅”的场景，天天在发生。法庭上法官要判“这起事故里，各方各占几成责任”。流行病学要判“这种病和那种暴露各有多大关系”。你复盘一次项目失败，也要判“哪个决策是主因”。这些事看起来无关，底层动作是同一个：把一堆搅在一起的因素，拆成“各占多少”。

这本书不给水晶球。它给的是一套拆解工具：当一堆因素共同导致了某个结果，怎么把“各占多少”定成一个可比的量。

最实在的例子就在法庭。法官不许给模糊答案，必须说出“这起事故70\%是他的责任，30\%是天气”。逼着自己把“谁推了一把”算成一个数字——这件事法院天天干，干了上千年。本书要做的，就是把这套已经被真实世界验证过的方法，从“判案”推广到“预测”。

所以全书的目标始终是：让你对一件还没发生的事，也能给出有依据的判断，并说清“这个判断依赖哪几个因素”。这正是预测的最小单元——反事实推演：如果当初没有X，今天会怎样。

\paragraph*{逻辑衔接}
工具是什么、为谁而写，先说清了。但工具要装在哪？社会上的因果具体长什么样？下一节“新结构：从动机到循环”给出全书地图，让你随时知道自己在哪。

% ==================================================================
\section{新结构：从动机到循环}

\begin{readingnote}
你有没有想过，为什么有的书读完只记得零散观点，有的书读完却带走一张地图？\\
差别往往不在内容，而在结构——章节之间是并列摆放，还是彼此咬合？
\end{readingnote}

全书六段，每一段都长在上一段遗留的问题上，不是随便摆的。

第一站，动机与目标（Part I）。要预测未来，先得看懂人想要什么。动机为什么是预测的必要输入？目标又是什么形状？

但是，为什么光有动机不够？因为动机到结果之间隔着环境和选择。所以第二站，预测与因果（Part II）。看懂了动机，凭什么能预测？答案：能预测就等价于有因果。相关不等于原因。

但是，一堆原因搅在一起怎么算？第三站，法庭上的时光机（Part III）。法律界几千年前就逼出了这套方法——必须给一个确定的数字。法律不是例子，是方法的证明。

但是，社会上的因果具体长什么样？第四站，社会的四大齿轮（Part IV）。拆成观念的病毒、经济、权力、法律四个通道，逐层加约束。文化是大集合，经济是加稀缺约束的子集，政治是加身份偏好的子子集，法律给三者的博弈结果提供定责框架。

最后，回到循环（Part V）。预测出的未来，会反过来塑造下一轮的动机。循环闭合。

\begin{figure}[h]
\centering
\begin{tikzpicture}[
node distance=0.7cm and 1.6cm,
every node/.style={draw, rectangle, rounded corners, align=center, font=\small}
]
\node (p0) {导论\\算得最准，预测最离谱};
\node (p1) [below=of p0] {Part I\\动机与目标};
\node (p2) [below=of p1] {Part II\\预测与因果};
\node (p3) [below=of p2] {Part III\\法庭上的时光机};
\node (p4) [below=of p3] {Part IV\\社会的四大齿轮};
\node (p5) [below=of p4] {Part V\\回到循环};
\draw[-{Stealth}] (p0) -- (p1);
\draw[-{Stealth}] (p1) -- (p2);
\draw[-{Stealth}] (p2) -- (p3);
\draw[-{Stealth}] (p3) -- (p4);
\draw[-{Stealth}] (p4) -- (p5);
\draw[-{Stealth}, dashed] (p5.east) .. controls +(2.2,0) and +(2.2,0) .. node[right, draw=none, align=left]{循环：预测\\塑造新动机} (p1.east);
\end{tikzpicture}
\caption{全书六段式地图。因果（Part II）从“想预测”里自然长出；法律（Part III）是这套方法最早的成熟应用现场；四大齿轮（Part IV）把同一套工具装到社会运转上。}
\label{fig:intro-dag}
\end{figure}

\paragraph*{逻辑衔接}
结构图有了。但一条环环相扣的链，会不会把读者锁死在单一读法里？下一节给出阅读策略，以及全书唯一的检验标准。

% ==================================================================
\section{如何阅读本书与检验标准}

\begin{readingnote}
你有没有遇到过：有些书结构清晰却枯燥，有些书生动却杂乱？\\
如果一本书想既保持逻辑统一、又允许灵活阅读——可能吗？本书试试。
\end{readingnote}

线性读法：第一次读，建议按顺序。这条链的依赖是真实的。不先有动机与目标，因果就悬空。不先有因果工具，法庭和方法应用就接不上。

跳跃读法：如果你已有背景，可按问题跳读。每一章开头的逻辑地图标明位置与前后依赖。每一节末尾的逻辑衔接显式写出“下一节将解决什么”。想直接看“各占多少”的读者，可以从Part III进入，缺什么再回Part II补。

正文用自然语言推进直觉，关键概念配一个“数学建模”盒子放形式化，可跳过，但跳过之后你要接受“知其然”的折扣。Part I和Part II零公式，公式从Part III起，且前面必有真实案例。

目标读者有三类：想理解社会系统如何被前向推导的自然科学背景者；愿意接受概率工具但不希望被公式淹没的社会科学研究者；对“社会能不能预测”本身感兴趣的普通读者。第三类请注意：本书的“预测”是概率断言，不是预言。说“降水概率70\%”的气象学家，和说“明天必然下雨”的巫师，做的是两种完全不同的事。

最后，把“这本书成功了”说清楚，否则“预测”会滑向它最俗的意义。本书的验收标准只有一句话：

\begin{quote}
\textbf{如果一个陌生人读完这本书，他能用书里的工具，对一个他没见过的社会事件给出概率判断，并说出“这个判断依赖于哪几个因素”——那么这本书就成功了。}
\end{quote}

相应地，“我们能说什么、我们不能说什么”必须分开写。本书能给的：对一个事件，一组因素的贡献排序与概率断言。对一个“如果当初没有X”的问题，一条可辩护的反事实推演。本书不能给的：明天的股价、下一位候选人、战争的具体日期。

\paragraph*{逻辑衔接}
标准立好了。接下来正式进入正文——从最日常的一句话开始：“他这么做，是因为……”。第1章“动机：解释的最小起点”将拆开这句话，看里面藏着什么假设。

\begin{thinkbox}
\textbf{想一想：}
\begin{enumerate}
  \item 2008年金融危机里，风险模型错把“房价一起涨”这个相关性当成了安全信号。你身边有没有把相关性当因果用的例子？如果要区分，你需要什么额外信息？
  \item 反事实练习：如果2007年评级机构没有给次级债贴上AAA，2008年的链条会怎么变？哪一环的概率会先掉？你会怎么估“各占多少”？
  \item 天气预报说“明天降水概率70\%”，和算命先生说“明天必有大雨”，两者在“可检验性”上有什么本质区别？本书要做的是哪一种？
\end{enumerate}
\end{thinkbox}

\end{document}
